
\begin{figure}[H]
    \centering
    \begin{tikzpicture}[
            scale=0.9,
            every node/.style={transform shape},
            bloque/.style={draw, minimum height=0.9cm, font=\footnotesize, align=center, inner sep=3pt},
            flecha/.style={->, thick}
        ]

        % Nodo de entrenamiento (superior)
        \node[bloque, minimum width=2.2cm, fill=blue!10] (train) at (0,0)
        {80\%\\Ajuste};
        % Nodo de validación interna, justo debajo
        \node[bloque, minimum width=1.3cm, fill=orange!15, below=1.0cm of train] (valint)
        {20\%\\Val.};

        % Etiqueta de etapa temporal arriba, centrada respecto ambos nodos
        \node[font=\footnotesize, above=0.42cm of train]
        {\textbf{Entrenamiento 20--40\%}};

        % Selección de hiperparámetros a la derecha, alineada verticalmente con train
        \node[bloque, minimum width=1.7cm, fill=gray!10, right=2.0cm of train] (select)
        {Selec.\\Hiperparam.};
        % Futuro a la derecha del anterior
        \node[bloque, minimum width=1.7cm, fill=green!12, right=1.7cm of select] (future)
        {Futuro\\40--60\%};

        % Flechas de proceso
        \draw[flecha] (train.east) -- (select.west);
        \draw[flecha] (valint.east) -- ([yshift=0.0cm]select.south west);

        % Unión de validación interna hacia la selección de hiperparámetros (desde la derecha de valint hacia abajo del bloque de selección)
        % Flecha de proceso hacia el futuro
        \draw[flecha] (select.east) -- (future.west);

    \end{tikzpicture}
    \caption[Esquema de validación interna para el ajuste de hiperparámetros]{Esquema de validación interna para el ajuste de hiperparámetros con los conjuntos de ajuste y validación dispuestos en paralelo.}
    \label{fig:ajuste_hiperparametros_validacion_interna}
\end{figure}

% Ajustes recomendados: scale ~0.80, fuente \footnotesize, minimum width y height menores, distancias proporcionadas.